\documentclass[11pt,a4paper]{article}
\usepackage[T1]{fontenc}
\usepackage{amsmath,amssymb,booktabs,geometry}
\geometry{margin=28mm}
\title{EPPA--III validation log: generalised clique covers}
\author{Independent audit and Lean formalisation}
\date{9 October 2026}
\begin{document}
\maketitle

\paragraph{Source and editing policy.}
This log refers to the supplied manuscript source \texttt{main-new(6).tex},
in particular Definitions 1.5--1.6, Lemma 3.7 and Claims 3.8--3.11.
The manuscript is frozen: corrections are proposed here and NOT applied
to its TeX. A successful Lean CI run does not certify any statement that
has not actually been expressed and proved in Lean.

\paragraph{Status.}
The initial Lean module models the graph given by Definition 1.6 and
proves its elementary induced-copy and adjacency properties.
The full classification Lemma 3.7 is \textbf{not yet validated}.
A separate exhaustive finite computation is complete for the
exceptional \(s=t=2\) matrix family.

\section*{Issues}

\paragraph{GCC--01 (confirmed): Taylor-double identification.}
The passage following Definition 1.6 says that, for \(G=K_t\),
diagonal connection \(X=M\) gives exactly the Taylor double of \(K_t\).
Definition 2.2 instead gives \(2K_t\), so the Taylor double corresponds
to \(X=0\). The matching connection has \(t\) additional cross-edges.
The analogous complement wording should also be revised.
Lean target: \texttt{matching\_not\_taylor\_clique}.

\paragraph{GCC--02 (proof repair required): Claim 3.8.}
The argument with coloured cycles needs a canonical connection profile,
not a profile dependent on arbitrary orderings of vertices inside the
\(D_j\)'s. One possible repair is to use, for every \(D_j\), the vector
whose \(i\)-th coordinate is the multiset of neighbourhoods of the
vertices of \(D_j\) in \(C_i\). If an automorphism extending a
transposition in \(C_i\) maps \(D_j\) to \(D_{j'}\), the coordinates other
than \(i\) agree. Unequal degrees of the transposed vertices force
the \(i\)-th coordinates to differ. Then the no-rainbow-cycle argument
must be stated for these actual profiles. This is a proposed repair
and has not been checked in Lean.

\paragraph{GCC--03 (serious gap): Claim 3.9.}
An extension of an arbitrary permutation of \(C_i\) may move the clique
\(D_j\). Thus the assertion that every \(k\)-subset of \(C_i\) occurs
as the neighbourhood of a vertex in this \emph{fixed} \(D_j\)
does not follow. Consequently, \(\binom tk\leq t\) and the deduction
\(k\in\{1,t-1\}\) are unproved. Counting neighbourhoods
across \emph{all} \(s\) cliques only gives
\(\binom tk\leq st\); this is generally insufficient. A genuine
stabiliser analysis is needed.

\paragraph{GCC--04 (gap): Claim 3.10.}
The colour-count argument on the clique-adjacency graph requires a
precisely identified action on the \(D_j\)'s. Its assertion of a
common number \(k_m\) of colour-\(m\) neighbours for every \(D_j\)
does not follow from transitivity on the \(C_i\)'s alone.
Even when all bipartite entries are isomorphic matchings,
independently reordering each \(D_j\) to turn every entry into
the same literal matching requires compatibility of the relative
permutations. A better strategy is to rule out off-diagonal
matching cases first and then normalise only the diagonal matchings.

\paragraph{GCC--05 (confirmed flaw): Claim 3.11 when \(s\geq3,t=2\).}
The calculation of \(N(I)\cap(H\setminus G)\) omits diagonal
connections. In particular, if \(X=1\), a transversal \(I\)
choosing one vertex of each \(C_j\) is adjacent to every vertex
of \(D_j\), for every \(j\). Hence
\(N(I)\cap(H\setminus G)=H\setminus G\) independently of \(I\).
The proof's claimed variation of this neighbourhood size therefore
fails without first excluding \(X=1\). This is an objection to
the \emph{argument}, not a counterexample to Lemma 3.7.
Lean target: \texttt{full\_diag\_hits\_every\_right}.

\paragraph{GCC--06 (gaps): remaining cases of Claim 3.11.}
In the \(s\geq3,t\geq3\) case, an extension fixing \(C_1\)
pointwise may permute the \(D_j\)'s; the proof assumes \(D_3\)
is fixed setwise. In the \(s=2,t\geq3\) case, the sentence
claiming that \emph{both} diagonal and off-diagonal entries
are matchings is not a consequence of the preceding claim:
the diagonal entries may still be \(0\) or \(1\).
These possibilities must be handled.

\paragraph{GCC--07 (confirmed, independent finite check): Claim 3.11.}
The claim that the Hadamard/Wagner pattern is the only
\emph{transitive} graph in the final \(s=t=2\) case is false.
Enumerating the 64 choices of diagonal entries from
\(\{0,1,M,\overline M\}\) and off-diagonal entries from
\(\{M,\overline M\}\) yields 24 transitive matrices in four
isomorphism classes. Exactly 16 matrices are EPPA witnesses
for the embedded \(2K_2\). The classes are:
\begin{center}
\begin{tabular}{@{}lrrr@{}}
\toprule
Representative & Transitive matrices & EPPA matrices & Edges\\
\midrule
\(\left(\begin{smallmatrix}0&M\\M&0\end{smallmatrix}\right)\) &4&4&8\\
\(\left(\begin{smallmatrix}1&M\\M&1\end{smallmatrix}\right)\) &4&4&16\\
\(\left(\begin{smallmatrix}M&M\\M&M\end{smallmatrix}\right)\) &8&0&12\\
\(\left(\begin{smallmatrix}M&M\\\overline M&M\end{smallmatrix}\right)\) &8&8&12\\
\bottomrule
\end{tabular}
\end{center}
The first two are ordinary covers after reindexing, the third
is the cube (fails EPPA for the designated copy), and the fourth
is the Wagner graph. The reproducible exhaustive check is
\texttt{checks/cover\_2k2.py}, written using Python's standard library.

\section*{Next formalisation goals}
\begin{enumerate}
\item Define partial automorphisms of the distinguished induced \(sK_t\)
and the exact EPPA condition for the inclusion in the cover graph.
\item Formalise the action on \(D\)-cliques of extensions of
\(\operatorname{Aut}(sK_t)\); validate the profile-cycle proof of Claim 3.8.
\item Find a correct stabiliser proof or a counterexample for Claim 3.9.
\item Establish the permitted matrix patterns, then reduce the small
exceptional case to the exact finite computation.
\end{enumerate}

\paragraph{Lean audit boundary.}
Only declarations explicitly listed in \texttt{CheckAxioms.lean}
are considered verified when the GitHub Actions build and axiom audit
are successful. Claims 3.8--3.11 and Lemma 3.7 are not yet on that list.
\end{document}
